\section{A formal source algebra in rank two}
\label{sec:vsrc}

The rank-two discussion of \Cref{sec:higher-rank-no-go} suggests separating
two questions: what algebraic shape a rank-$r$ source should have, and
whether actual cycle classes realize that shape.  This section treats the
first question with the simplest possible device, the exterior algebra on
$r$ formal symbols.  The top product $J_1\cdots J_r$ plays the role of a
determinant: it is nonzero exactly when the symbols are independent.  We
call this algebra the \emph{virtual source algebra} $\vsrc_r$; the word
``virtual'' signals that the symbols are not, by themselves, cycle classes.

\begin{definition}[Virtual source algebra]\label{def:apparatus:vsrc-algebra}
Let $V_r=\Q J_1\oplus\cdots\oplus\Q J_r$ and $\vsrc_r=\Lambda_{\Q}(V_r)$,
the exterior algebra.  In rank $2$ it has ordered basis
$1,\ \Jone,\ \Jtwo,\ \Jtwelve:=\Jone\Jtwo$ and relations
\[
  \Jone^2=0,\qquad \Jtwo^2=0,\qquad
  \Jone\Jtwo=\Jtwelve,\qquad
  \Jtwo\Jone=-\Jtwelve.
\]
The element $\Jtwelve$ is the determinant-shaped witness of independence
of $\Jone$ and $\Jtwo$.%
\footnote{Lean constructs the rank-two multiplication table on integer
coordinates and proves associativity, the unit laws and additive and
scalar bilinearity.  Rational coefficients, other ranks and the dimension
count are treated here on paper.}
\end{definition}

\Cref{fig:vsrc-wedge} displays the rank-two algebra.
\begin{figure}[tbp]
\centering
\begin{tikzpicture}[
  font=\small,
  node/.style={draw, rounded corners=2pt, align=center, minimum width=1.6cm,
    minimum height=0.8cm, fill=black!3},
  arrow/.style={-{Latex[length=2mm]}, thick}
]
\node[node] (one) at (0,2.2) {\(1\)};
\node[node] (j1) at (-2.2,0.7) {\(\Jone\)};
\node[node] (j2) at (2.2,0.7) {\(\Jtwo\)};
\node[node] (j12) at (0,-1.1) {\(\Jtwelve\)};

\draw[arrow] (j1) -- node[above, sloped] {\(\wedge\,\Jtwo\)} (j12);
\draw[arrow] (j2) -- node[above, sloped] {\(\wedge\,\Jone=-\)} (j12);
\draw[arrow] (one) -- node[left] {\(\wedge\,\Jone\)} (j1);
\draw[arrow] (one) -- node[right] {\(\wedge\,\Jtwo\)} (j2);
\draw[loop left, looseness=6, -{Latex[length=2mm]}] (j1) to node[left]
  {\(\Jone^2=0\)} (j1);
\draw[loop right, looseness=6, -{Latex[length=2mm]}] (j2) to node[right]
  {\(\Jtwo^2=0\)} (j2);
\node[align=center, font=\footnotesize] at (0,-2.25)
  {\(\dim_{\Q}\vsrc_2=4\), with \(\Jone\Jtwo=\Jtwelve\) the determinant-shaped witness};
\end{tikzpicture}
\caption{The rank-\(2\) \(\vsrc\) exterior-algebra calibration.}
\label{fig:vsrc-wedge}
\end{figure}


\begin{theorem}[Stage I: consistency]\label{thm:apparatus:vsrc-stage-i-consistency}
$\vsrc_2$ is an associative, unital, graded-commutative $\Q$-algebra.
Every element has a unique normal form
$a_{\varnothing}1+a_1\Jone+a_2\Jtwo+a_{12}\Jtwelve$ with $a_\bullet\in\Q$,
so $\dim_{\Q}\vsrc_2=4$.  In lower rank, $\vsrc_0=\Q$ and
$\vsrc_1=\Q\oplus\Q J_1$ with $J_1^2=0$.%
\footnote{Lean proves the unit laws, the four relations and the
distinctness of the basis symbols for the integer-coordinate model, and
records the numbers $2^0,2^1,2^2$; it does not prove the dimension
statement over $\Q$.}
\end{theorem}

\begin{proof}
Take the $\Q$-vector space with basis indexed by the subsets of
$\{1,2\}$.  Define the product of two basis vectors to be $0$ if the
subsets meet, and otherwise to be the basis vector of the union times the
sign of the permutation that sorts the concatenated index list; extend
bilinearly.  This gives the displayed relations, and the empty subset is a
unit.  Associativity holds on basis vectors because both
parenthesizations vanish when an index repeats, and otherwise both
produce the same union with the sign of the same sorting permutation.
Moving a degree-$a$ monomial past a degree-$b$ monomial costs
$(-1)^{ab}$, which is graded commutativity.  The four subsets give a basis,
so the dimension is $4$; the same construction with $\{1\}$ or
$\varnothing$ gives $\vsrc_1$ and $\vsrc_0$.
\end{proof}

\begin{remark}[No rank-$r$ realization is supplied here]\label{rmk:apparatus:vsrc-no-licensed-descent}
For $r\ge2$, nothing in this paper provides a map
\[
  \vsrc_r(E)\longrightarrow \mathrm{CH}^*(\mathcal X)_{\Q},
\]
or into a Selmer or Mordell--Weil group, that realizes $J_1,\ldots,J_r$ as
independent cycle classes.  The Gram-matrix statement
\Cref{thm:apparatus:gram-matrix-shadow} shows that heights alone do not
supply such provenance.  Constructions of cycle classes with arithmetic
meaning, such as Kato's Euler system~\citep{Kato2004} and generalized
Heegner cycles~\citep{BertoliniDarmonPrasanna2013}, come with their own
hypotheses; we do not use them, and we make no claim about whether they
realize $\vsrc_r$ for a given curve.
\end{remark}

In rank $1$ the situation is classical.  Under the Gross--Zagier
hypotheses~\citep{GrossZagier1986,Kolyvagin1990} for an imaginary
quadratic field $K$, there is a Heegner point $P_K\in E(K)$, and one can
send $J_1$ to $P_K$.

\begin{theorem}[Stage II: rank one]\label{thm:apparatus:vsrc-stage-ii-gz-partial}
Send $J_1\mapsto P_K$ and the formal height $h_{\vsrc}(J_1,J_1)$ to the
N\'eron--Tate height $\langle P_K,P_K\rangle_{\mathrm{NT}}$.
\begin{enumerate}
\item The $1\times1$ Gram determinant of the image source is
  $\langle P_K,P_K\rangle_{\mathrm{NT}}$.  If $P_K$ has infinite order and
  $E(K)$ has rank $1$, this equals $[\Lambda:\Z P_K]^2\Reg(E/K)$, where
  $\Lambda=E(K)/E(K)_{\tors}$ and $[\Lambda:\Z P_K]$ is the index of the
  image of $\Z P_K$.
\item The Gross--Zagier identity relating $L'(E/K,1)$ to
  $\langle P_K,P_K\rangle_{\mathrm{NT}}$ is not determined by these
  symbolic data: two formal data sets can share $P_K$ and
  $\langle P_K,P_K\rangle_{\mathrm{NT}}$ while having different analytic
  quotients.  Relative to the symbolic calculus the analytic identity is an
  additional input, which we record by a nonzero residual
  $\Xi_{\mathrm{GZ}}$.\footnote{Lean checks the symbolic assignments $J_1\mapsto P_K$ and
$h_{\vsrc}(J_1,J_1)\mapsto$ (N\'eron--Tate height), and records the
declared residual value $\Xi_{\mathrm{GZ}}=1$; it constructs neither the
Heegner point nor any height.}
\end{enumerate}
\end{theorem}

\begin{proof}
(1) The Gram determinant of a single vector is its height.  If $Q$
generates $\Lambda$ and $P_K\equiv mQ$ modulo torsion, then
$\langle P_K,P_K\rangle_{\mathrm{NT}}=m^2\langle Q,Q\rangle_{\mathrm{NT}}
=m^2\Reg(E/K)$ with $m=[\Lambda:\Z P_K]$.  (2) The symbolic data contain
neither $L'(E/K,1)$ nor the period, so they cannot constrain the analytic
side; the identity is the content of the Gross--Zagier theorem.
\end{proof}

In rank $2$ the corresponding statement is a one-way translation.  For
$389a1$ the Tamagawa product and the torsion order are $1$; we write
\[
  \Xi_{\mathrm{Sc}}
  =A_E-\Omega_E^+\,\Reg(E/\Q)\,\#\Sha(E/\Q)
\]
for the scalar BSD residual, assuming $\Sha(E/\Q)$ finite, which is not
known for this curve.  A \emph{rank-two realization} consists of a
$\Q$-vector space $W$ of classes (for instance cycle classes modulo
homological equivalence), a linear map $f:V_2\to W$, a symmetric
bilinear height pairing on $W$, and a bridge
$\mathcal B^{(2)}_{\mathrm{an\text{-}ht}}$ identifying $A_E$ with
$\Omega_E^+\#\Sha$ times the Gram determinant of $(f(\Jone),f(\Jtwo))$
for that pairing.  Put $r_2=\Lambda f:\vsrc_2\to\Lambda_{\Q}(W)$; then
$z_{12}=r_2(\Jtwelve)=f(\Jone)\wedge f(\Jtwo)$, and $z_{12}\neq0$ exactly
when $f(\Jone),f(\Jtwo)$ are linearly independent.  Let
$\Xi_{\mathrm{Mot}}$ be a \emph{motivic residual}, which by definition
vanishes exactly when a rank-two realization with $z_{12}\neq0$ is given.
All of these data are supplied; none is constructed here.

\begin{theorem}[Stage III: one-way translation for $389a1$]\label{thm:apparatus:vsrc-stage-iii-translation}
For the fixed pair $(\Xi_{\mathrm{Mot}},\Xi_{\mathrm{Sc}})$ just defined,
suppose given the analytic--height bridge in the form of the implication
$\Xi_{\mathrm{Mot}}=0\Rightarrow\Xi_{\mathrm{Sc}}=0$.  Then a rank-two
realization yields the scalar identity
$A_E=\Omega_E^+\Reg(E/\Q)\#\Sha(E/\Q)$.  The reverse implication
$\Xi_{\mathrm{Sc}}=0\Rightarrow\Xi_{\mathrm{Mot}}=0$ is not provided by
the scalar data; it holds only if supplied as separate content.%
\footnote{The formalization records this as a conditional schema: for a
fixed pair of residuals, the forward implication is projected from the
supplied bridge, and the reverse from supplied external content.  Lean
does not construct a realization or a bridge.}
\end{theorem}

\begin{proof}
The forward statement is the supplied implication.  No reverse construction of a realization or
of an independence certificate from the scalar equality is supplied
here.  (\Cref{thm:apparatus:gram-matrix-shadow} establishes only that
provenance data cannot be recovered from the Gram matrix in its record
model; it does not exclude a separately defined construction.)
\end{proof}

The statement must be read for the fixed pair.  A universal version,
asserting $\Xi_{\mathrm{Mot}}=0\Rightarrow\Xi_{\mathrm{Sc}}=0$ for all
values of the two residuals, would be false, for instance at
$(\Xi_{\mathrm{Mot}},\Xi_{\mathrm{Sc}})=(0,1)$.

\begin{remark}[The source reading carries more data]\label{rmk:apparatus:vsrc-rank-r-stronger}
The scalar leading-coefficient formula records a number.  A rank-$r$
source statement in the sense above also asks for a nonzero independent
wedge $J_1\wedge\cdots\wedge J_r$ realized by actual classes.  In this
sense the source reading carries more data than the scalar equality.  We
do not claim a general theorem that scalar BSD is strictly weaker than any
such source statement; the comparison above is for the fixed pair and the
supplied bridge.
\end{remark}
